% Fiche C — recto-verso A4 paysage
% Socle commun C / C++ : tout ce qui est partagé (types, opérateurs, contrôle, pointeurs,
% mémoire, préprocesseur, compilation) est ici ; la fiche C++ ne contient que ce qu'ajoute le C++.
% Compiler : make (ici ou à la racine), ou xelatex c_cheat.tex — une seule passe suffit
% Apparence : ../librairie (fiche.sty + theme-c.tex). Macros de contenu :
%   colonne[titre][couleur] (3 par page, séparées par \colbreak) · bloc{titre}{cat}
%   catégories : syn typ ctl ex io lib (couleurs et icônes dans le thème)
%   code · \V|...| code en ligne · «...» dans le code = à remplir
%   \sub sous-titre · \note remarque · \piege piège · \badge{TXT} · schema
%   optab étiquette | texte · tab3 étiquette | code | remarque · compare · grille{n}
%   (largeurs mesurées par la librairie : celles écrites entre [] et {} sont ignorées)
% Le code est en police à chasse fixe : les commentaires sont alignés à la main
% (attention : les « » ne sont pas affichés, ils ne comptent pas dans l'alignement).
% Après compilation : grep "\[fiche\]" c_cheat.log (place libre, collisions).
\documentclass{article}
\usepackage[c]{../librairie/fiche}   % apparence : ../librairie (fiche.sty + theme-c.tex)
% --- Langage C pour listings ---
% Classes de mots-clés : [1] rose · [2] cyan · [3] vert · [4] couleur du thème
% (styles dans ../librairie/fiche.sty). Chargé par c_cheat.tex.
\lstdefinelanguage{CODE}{
  sensitive=true, alsoletter={\#},
  morekeywords=[1]{if,else,for,while,do,return,break,continue,switch,case,default,
    goto,sizeof,struct,union,enum,typedef,static,extern,const,volatile,inline,
    register,auto,restrict},
  morekeywords=[2]{int,float,double,char,void,short,long,signed,unsigned,bool,_Bool,
    size_t,ptrdiff_t,int8_t,uint8_t,int16_t,uint16_t,int32_t,uint32_t,int64_t,
    uint64_t,FILE,node_t,binop_t},
  morekeywords=[3]{printf,scanf,fprintf,fscanf,snprintf,puts,putchar,getchar,fopen,
    fclose,fgets,fputs,fread,fwrite,perror,malloc,calloc,realloc,free,exit,qsort,
    strlen,strcpy,strncpy,strcmp,strcat,memcpy,memset,NULL,EOF,stdin,stdout,stderr,
    \#include,\#define,\#ifdef,\#ifndef,\#if,\#else,\#endif,\#undef,\#pragma,
    __FILE__,__LINE__,__func__,__VA_ARGS__},
  morecomment=[l]{//}, morecomment=[s]{/*}{*/}, morestring=[b]", morestring=[b]',
}
   % coloration syntaxique : \lstdefinelanguage{CODE}

\begin{document}
%%%%%%%%%%%%%%%%%%%%%%%%%%%%%%%%%%%%%%%%%%%%%%%%%%%%%%%%%%%%%%%%%%%%%%%%%%%%%%
%                         RECTO : LE LANGAGE
%%%%%%%%%%%%%%%%%%%%%%%%%%%%%%%%%%%%%%%%%%%%%%%%%%%%%%%%%%%%%%%%%%%%%%%%%%%%%%
\entete{Langage C}{recto · le langage}
\noindent
% ======================= colonne 1 : types et opérateurs =======================
\begin{colonne}[types · opérateurs][col-a]
\begin{bloc}{Types de base}{typ}
\begin{tab3}[14mm]{27mm}
Entiers     & \V|char short int|    & 1 · 2 · 4 octets (typique) \\
            & \V|long|              & 8 (Linux 64) · 4 (Windows) \\
            & \V|long long|         & $\geq$ 8 octets \\
            & \V|unsigned int|      & $\geq 0$, calcul modulo $2^n$ \\
Fixes       & \V|int8_t uint32_t|   & taille exacte · \V|<stdint.h>| \\
Réels       & \V|float double|      & 4 · 8 octets (IEEE 754) \\
Booléen     & \V|bool true false|   & \V|<stdbool.h>| (mots-clés en C23) · 0 = faux \\
Tailles     & \V|size_t|            & non signé, type de \V|sizeof| \\
Vide        & \V|void|              & pas de valeur · \V|void *| générique
\end{tab3}
\sub{Littéraux}
\begin{code}
42   052   0x2A   0b101010    // déc, oct, hexa, bin (C23)
42u  42L   42UL   42LL        // suffixes de type
3.14  3.14f  1e-3             // double, float, 0.001
'A'  '\n'  '\0'  '\x41'       // caractères (= entiers)
"texte"                       // char[] terminé par '\0'
\end{code}
\sub{Conversions}
\begin{code}
7 / 2           // = 3   (division entière)
7 / 2.0         // = 3.5 (un réel suffit)
(double)a / b   // cast avant la division
(int)3.9        // = 3   (troncature vers 0)
-7 % 3          // = -1  (signe du dividende)
-1 < 1u         // FAUX : -1 converti en unsigned
\end{code}
\note{Promotion : \V|char| et \V|short| deviennent \V|int| dans un calcul ;
mélange signé / non signé $\Rightarrow$ non signé.}
\end{bloc}

\begin{bloc}{Opérateurs par priorité}{syn}
\begin{tab3}[3mm]{31mm}
1  & \V|() [] -> . x++ x--|   & appel, indice, membre \\
2  & \V|! ~ ++x -x *p &x|     & unaires · droite $\to$ gauche \\
   & \V|(type) sizeof|        & cast, taille en octets \\
3  & \V|* / %|                & \\
4  & \V|+ -|                  & \\
5  & \V|<< >>|                & décalages \\
6  & \V|< <= > >=|            & \\
7  & \V|== !=|                & \\
8  & \V|&|                    & ET bit à bit \\
9  & \V|^|                    & XOR bit à bit \\
10 & \V/|/                    & OU bit à bit \\
11 & \V|&&|                   & ET logique (court-circuit) \\
12 & \V/||/                   & OU logique (court-circuit) \\
13 & \V|c ? a : b|            & condition · droite $\to$ gauche \\
14 & \V|= += -= *= <<= &=|    & affectations · droite $\to$ gauche \\
15 & \V|a, b|                 & évalue a puis b, vaut b
\end{tab3}
\piege{\V|x & MASK == 0| se lit \V|x & (MASK == 0)| : parenthéser les opérations bit à bit.}
\end{bloc}

\begin{bloc}{Pièges classiques}{syn}
\begin{optab}[2mm]
· & \V|if (x = 0)| affecte : écrire \V|if (x == 0)| \\
· & \V|;| après \V|if (...)| ou \V|for (...)| : corps vide \\
· & \V|switch| sans \V|break| : enchaîne les cas \\
· & dépassement d'un \V|int| signé : comportement indéfini \\
· & \V|scanf("%d", x)| : oubli du \V|&| \\
· & chaîne sans place pour le \V|'\0'| final
\end{optab}
\end{bloc}
\end{colonne}\colbreak
% ======================= colonne 2 : contrôle et fonctions =======================
\begin{colonne}[contrôle · fonctions][col-b]
\begin{bloc}{Structures de contrôle}{ctl}
\begin{code}
if («cond») {                 // cond : 0 = faux, sinon vrai
    «instructions»
} else if («cond2») {
    «instructions»
} else {
    «instructions»
}

switch («expr_entière») {
    case 1:  «instr»; break;
    case 2:                   // sans break : continue
    case 3:  «instr»; break;  //   dans le case suivant
    default: «instr»;
}

for (int i = 0; i < n; i++) { ... }
while («cond») { ... }
do { ... } while («cond»);    // au moins 1 passage
break;                        // sort de la boucle
continue;                     // passe au tour suivant
goto fin;  ...  fin: ;        // sortie d'erreur unique
\end{code}
\end{bloc}

\begin{bloc}{Fonctions}{ctl}
\begin{code}
int add(int a, int b);        // prototype (avant usage)
int add(int a, int b) {       // définition
    return a + b;
}
void swap(int *a, int *b) {   // modifier l'appelant :
    int t = *a;               //   passer l'adresse
    *a = *b;
    *b = t;
}
swap(&x, &y);
void f(void);                 // aucun paramètre
double moy(const int *t, size_t n);  // tableau + taille
static int aide(int x);       // privée au fichier .c
int main(int argc, char *argv[]) {   // argv[0] = nom
    return 0;                 // 0 = succès
}
\end{code}
\note{Passage \textbf{toujours par valeur} : la fonction reçoit une copie.
Un tableau passé en paramètre devient un pointeur sur son 1\textsuperscript{er} élément
(sa taille est perdue).}
\end{bloc}

\begin{bloc}{Portée et durée de vie}{ctl}
\begin{optab}[19mm]
Locale       & pile, détruite en fin de bloc, \textbf{non initialisée} \\
Globale      & toute la durée du programme, initialisée à 0 \\
\V|static| loc. & garde sa valeur entre deux appels \\
\V|static| glob. & visible dans ce fichier seulement \\
\V|extern|   & déclare une globale définie ailleurs \\
\V|const|    & lecture seule \\
\V|volatile| & relue en mémoire à chaque accès (registre, interruption)
\end{optab}
\begin{code}
int compteur(void) {
    static int n = 0;         // initialisé une seule fois
    return ++n;               // 1, 2, 3...
}
\end{code}
\end{bloc}
\end{colonne}\colbreak
% ======================= colonne 3 : pointeurs, tableaux, chaînes =======================
\begin{colonne}[pointeurs · chaînes][col-c]
\begin{bloc}{Pointeurs}{typ}
\begin{code}
int x = 5;
int *p = &x;                  // p = adresse de x
*p = 7;                       // x vaut 7
int **pp = &p;                // pointeur de pointeur
p = NULL;                     // ne pointe sur rien
void *v = &x;                 // générique (pas de *v)
p + 1                         // avance de sizeof(int)
p[i]  ==  *(p + i)            // indice = arithmétique
q - p                         // écart en éléments
s->champ  ==  (*s).champ      // membre via pointeur
\end{code}
\sub{Lire une déclaration}
\begin{tab3}[0mm]{24mm}
 & \V|const int *p|   & donnée constante, p modifiable \\
 & \V|int *const p|   & p constant, donnée modifiable \\
 & \V|int *t[10]|     & tableau de 10 pointeurs \\
 & \V|int (*t)[10]|   & pointeur sur un tableau de 10 \\
 & \V|int *f(int)|    & fonction qui renvoie \V|int *| \\
 & \V|int (*f)(int)|  & pointeur de fonction
\end{tab3}
\sub{Pointeurs de fonction}
\begin{code}
typedef int (*binop_t)(int, int);
binop_t op = add;             // nom = adresse
int r = op(2, 3);             // appel indirect
int cmp(const void *a, const void *b) {
    int x = *(const int *)a, y = *(const int *)b;
    return (x > y) - (x < y); // -1, 0 ou 1
}
qsort(t, n, sizeof t[0], cmp);   // tri générique
\end{code}
\end{bloc}

\begin{bloc}{Tableaux et chaînes}{typ}
\begin{code}
int t[5] = {1, 2};            // reste mis à 0
int u[] = {1, 2, 3};          // taille déduite : 3
int m[3][4];                  // m[i][j], ligne par ligne
size_t n = sizeof t / sizeof t[0];   // nb d'éléments
char s[] = "abc";             // 4 octets, '\0' compris
char *lit = "abc";            // littéral NON modifiable
void g(int m[][4], int rows); // 2D : colonnes fixées
\end{code}
\piege{\V|sizeof t| dans une fonction donne la taille du pointeur ; aucun contrôle des
bornes : indices de 0 à n$-$1.}
\end{bloc}

\begin{bloc}{Bibliothèque standard}{lib}
\begin{tab3}[15mm]{24mm}
<string.h> & \V|strlen(s)|        & longueur sans le \V|'\0'| \\
           & \V|strcpy(d, s)|     & copie (d assez grand) \\
           & \V|strncpy(d, s, n)| & \V|'\0'| non garanti \\
           & \V|strcat(d, s)|     & ajoute s à la fin de d \\
           & \V|strcmp(a, b)|     & $<0$, 0, $>0$ \\
           & \V|strchr strstr|    & $\to$ pointeur ou NULL \\
           & \V|memcpy memmove|   & copie n octets ; recouvrement : move \\
           & \V|memset(p, v, n)|  & n octets mis à v \\
<stdlib.h> & \V|strtol(s, &e, b)| & chaîne $\to$ long, base b \\
           & \V|rand() % n|       & entier de 0 à n$-$1 \\
           & \V|srand(time(NULL))| & graine, une seule fois \\
<math.h>   & \V|sqrt pow fabs|    & lier avec \V|-lm| \\
<ctype.h>  & \V|isdigit toupper|  & test d'un caractère \\
<assert.h> & \V|assert(p != NULL)| & arrête si faux
\end{tab3}
\piege{\V|a == b| compare les adresses, pas le texte : \V|strcmp(a, b) == 0|.}
\end{bloc}
\end{colonne}

%%%%%%%%%%%%%%%%%%%%%%%%%%%%%%%%%%%%%%%%%%%%%%%%%%%%%%%%%%%%%%%%%%%%%%%%%%%%%%
%                     VERSO : DONNÉES, MÉMOIRE ET OUTILS
%%%%%%%%%%%%%%%%%%%%%%%%%%%%%%%%%%%%%%%%%%%%%%%%%%%%%%%%%%%%%%%%%%%%%%%%%%%%%%
\newpage
\entete{Langage C}{verso · données, mémoire et outils}
\noindent
% ======================= colonne 4 : types composés, préprocesseur, bits =======================
\begin{colonne}[données · macros][col-a]
\begin{bloc}{struct, union, enum}{typ}
\begin{code}
struct point { int x, y; };
struct point a = {1, 2};
struct point b = {.y = 5};    // désignateur : x = 0
typedef struct node {
    int val;
    struct node *next;        // auto-référence
} node_t;
node_t n, *pn = &n;
n.val = 3;                    // accès direct
pn->val = 3;                  // via pointeur
b = a;                        // copie champ à champ
union mot {                   // champs superposés :
    uint32_t w;               //   même mémoire,
    uint8_t  o[4];            //   taille du plus grand
};
enum etat { IDLE, RUN = 5, STOP };   // 0, 5, 6
struct ctrl {                 // champs de bits
    unsigned en   : 1;
    unsigned mode : 3;
};
\end{code}
\note{\V|a == b| interdit entre structures : comparer champ par champ.}
\sub{Alignement et padding}
\begin{code}
struct s1 { char c; int i; char d; };  // sizeof = 12
struct s2 { int i; char c; char d; };  // sizeof = 8
\end{code}
\begin{schema}[x=1mm,y=1mm,font=\fontcode]
  \foreach \i in {1,2,3,9,10,11} \fill[gutter!35] (\i*3.4,0) rectangle +(3.4,3.4);
  \foreach \i/\c in {0/c,4/i,5/i,6/i,7/i,8/d} \node[fill=cat!20,minimum size=3.4mm,inner sep=0pt,anchor=south west] at (\i*3.4,0) {\c};
  \draw[trait,line width=0.4pt] (0,0) grid[step=3.4] (40.8,3.4);
  \node[anchor=west,text=note,font=\fontbloc] at (42,1.7) {s1 : 12 octets dont 6 de bourrage};
\end{schema}
\note{Chaque champ est aligné sur sa taille, la structure sur son plus grand champ :
ranger les champs du plus grand au plus petit.}
\end{bloc}

\begin{bloc}{Préprocesseur}{syn}
\begin{code}
#include <stdio.h>            // en-tête système
#include "module.h"           // en-tête du projet
#define N 100                 // constante (pas de ;)
#define SQ(x) ((x) * (x))     // tout parenthéser
#define STR(x) #x             // STR(a+b) -> "a+b"
#define CAT(a, b) a##b        // CAT(x, 1) -> x1
#ifdef DEBUG                  // gcc -DDEBUG
#define LOG(...) fprintf(stderr, __VA_ARGS__)
#else
#define LOG(...)              // disparaît
#endif
__FILE__ __LINE__ __func__    // position courante
\end{code}
\piege{\V|SQ(i++)| incrémente deux fois : une macro recopie le texte.}
\end{bloc}

\begin{bloc}{Bits et registres}{ex}
\begin{code}
x |=  (1u << n);              // met le bit n à 1
x &= ~(1u << n);              // met le bit n à 0
x ^=  (1u << n);              // inverse le bit n
if (x & (1u << n)) { ... }    // teste le bit n
v = (x >> 4) & 0xF;           // extrait les bits 7..4
#define GPIO_ODR (*(volatile uint32_t *)0x40020014)
GPIO_ODR |= 1u << 5;          // registre mappé
\end{code}
\end{bloc}
\end{colonne}\colbreak
% ======================= colonne 5 : mémoire =======================
\begin{colonne}[mémoire][col-b]
\begin{bloc}{Zones mémoire}{lib}
\begin{schema}[x=1mm,y=1mm,font=\fontcode]
  \foreach \x/\w/\n in {0/9/.text,9/12/.rodata,21/9/.data,30/8/.bss,38/10/tas}
    \node[draw=trait,fill=cat!20,minimum width=\w mm,minimum height=4.2mm,inner sep=0pt,anchor=west] at (\x,0) {\n};
  \node[draw=trait,fill=cat!20,minimum width=10mm,minimum height=4.2mm,inner sep=0pt,anchor=east] at (78,0) {pile};
  \draw[-{Latex[length=1.4mm]},draw=cat,line width=0.7pt] (49,0) -- (56,0);
  \draw[-{Latex[length=1.4mm]},draw=cat,line width=0.7pt] (67,0) -- (60,0);
  \draw[-{Latex[length=1.2mm]},draw=gutter] (0,-3.8) -- (78,-3.8)
    node[midway,fill=white,text=note,inner sep=1pt,font=\fontbloc] {adresses croissantes};
\end{schema}
\begin{optab}[13mm]
Pile    & locales, paramètres ; libérée en fin de bloc ; quelques Mo \\
Tas     & \V|malloc| / \V|free| : libéré à la main \\
.data   & globales et \V|static| initialisées \\
.bss    & globales et \V|static| non initialisées (mises à 0) \\
.rodata & constantes, littéraux \V|"..."| \\
.text   & code machine
\end{optab}
\end{bloc}

\begin{bloc}{Allocation dynamique}{lib}
\begin{code}
int *t = malloc(n * sizeof *t);   // non initialisé
if (t == NULL) {                  // toujours tester
    perror("malloc");
    exit(EXIT_FAILURE);
}
int *z = calloc(n, sizeof *z);    // mis à 0
int *tmp = realloc(t, 2 * n * sizeof *t);
if (tmp != NULL) t = tmp;         // échec : t intact
free(t);
t = NULL;                         // anti-pointeur pendant
\end{code}
\piege{un \V|free| par \V|malloc| ; jamais deux fois ; ne plus utiliser la zone après
\V|free| ; jamais sur une variable de pile.}
\sub{Matrice dynamique}
\begin{code}
int **m = malloc(rows * sizeof *m);
for (int i = 0; i < rows; i++)
    m[i] = malloc(cols * sizeof **m);
m[i][j] = 0;                      // accès comme en 2D
for (int i = 0; i < rows; i++)    // libérer dans
    free(m[i]);                   //   l'ordre inverse
free(m);
\end{code}
\sub{Pointeur pendant}
\begin{code}
int *mauvais(void) {
    int loc = 42;
    return &loc;              // loc détruite au retour !
}
\end{code}
\end{bloc}

\begin{bloc}{Exemple : liste chaînée}{ex}
\begin{code}
node_t *push(node_t *head, int v) {
    node_t *n = malloc(sizeof *n);
    if (n == NULL) return head;
    n->val  = v;
    n->next = head;
    return n;                     // nouvelle tête
}
node_t *l = NULL;                 // liste vide
l = push(l, 1);
l = push(l, 2);                   // l : 2 -> 1
for (node_t *p = l; p != NULL; p = p->next)
    printf("%d ", p->val);
while (l != NULL) {               // libération
    node_t *suiv = l->next;
    free(l);
    l = suiv;
}
\end{code}
\end{bloc}
\end{colonne}\colbreak
% ======================= colonne 6 : E/S, compilation =======================
\begin{colonne}[E/S · compilation][col-c]
\begin{bloc}{E/S <stdio.h>}{io}
\begin{tab3}[0mm]{22mm}
 & \V|%d %u %ld %lld| & int, unsigned, long, long long \\
 & \V|%zu %p|         & size\_t, pointeur \\
 & \V|%x %X %o|       & hexa, octal \\
 & \V|%f %e %g|       & double (float promu) \\
 & \V|%c %s %%|       & caractère, chaîne, \% \\
 & \V|%5d %-5d %05d|  & largeur, à gauche, zéros \\
 & \V|%.2f %8.3f|     & précision · \V|%#x| : préfixe \V|0x|
\end{tab3}
\begin{code}
printf("x=%d y=%.2f\n", x, y);
fprintf(stderr, "erreur %d\n", code);
snprintf(buf, sizeof buf, "%d", x);   // borné
int nb = scanf("%d %lf", &i, &d); // nb lus ; %lf : double
int c = getchar();            // int, car EOF = -1
\end{code}
\sub{Fichiers}
\begin{code}
FILE *f = fopen("data.txt", "r"); // "w" "a" "rb" "r+"
if (f == NULL) { perror("fopen"); return 1; }
char ligne[256];
while (fgets(ligne, sizeof ligne, f) != NULL)
    fputs(ligne, stdout);     // garde le '\n'
fprintf(f, "%d\n", x);        // écriture formatée
fread(buf, sizeof *buf, n, f);    // binaire
fclose(f);
\end{code}
\end{bloc}

\begin{bloc}{Compilation et débogage}{io}
\begin{tab3}[14mm]{12mm}
Préproc. & \V|gcc -E| & \V|#include|, macros $\to$ \V|.i| \\
Compil.  & \V|gcc -S| & C $\to$ assembleur \V|.s| \\
Assemb.  & \V|gcc -c| & $\to$ objet \V|.o| \\
Liens    & \V|gcc|    & \V|.o| + bibliothèques $\to$ exécutable
\end{tab3}
\begin{code}
gcc -Wall -Wextra -std=c17 -g -o prog main.c mod.c
gcc -c mod.c                  // -> mod.o seul
gcc main.o mod.o -lm -o prog  // + libm (math.h)
\end{code}
\begin{tab3}[0mm]{28mm}
 & \V|-O0 -O2|       & optimisation \\
 & \V|-I dir -L dir| & chemins des .h, des .a / .so \\
 & \V|-DNOM=val|     & \V|#define| depuis la ligne de commande \\
 & \V|-fsanitize=address| & débordements, fuites
\end{tab3}
\sub{module.h}
\begin{code}
#ifndef MODULE_H              // garde d'inclusion
#define MODULE_H
extern int compteur;          // déclaration seule
int api(int x);               // prototypes publics
#endif
\end{code}
\sub{Makefile}
\begin{code}
CFLAGS = -Wall -Wextra -std=c17 -g
prog: main.o mod.o            // cible: dépendances
    $(CC) $^ -lm -o $@        // TAB devant la commande
%.o: %.c mod.h                // règle générique
    $(CC) $(CFLAGS) -c $<     // $< 1re dépendance
\end{code}
\sub{gdb et valgrind}
\begin{tab3}{}
 & \texttt{gdb ./prog}            & programme compilé avec \V|-g| \\
 & \texttt{break f} · \texttt{run} & point d'arrêt dans f, lancer \\
 & \texttt{next} · \texttt{step}  & ligne suivante · entrer dans l'appel \\
 & \texttt{print x} · \texttt{bt} & valeur · pile d'appels \\
 & \texttt{valgrind ./prog}       & fuites : \texttt{-{}-leak-check=full}
\end{tab3}
\end{bloc}
\end{colonne}
\end{document}
